\textbf{结论名称：}蒙日圆（椭圆外两垂直切线的交点轨迹）

\textbf{适用条件：}标准椭圆 \( \dfrac{x^{2}}{a^{2}}+\dfrac{y^{2}}{b^{2}}=1\;(a>b>0) \)，椭圆外一点向椭圆引两条切线，且两条切线互相垂直。

\textbf{变量说明：} \( a,b \) 分别为椭圆的长半轴长和短半轴长 \( (a>b>0) \)；动点 \( P(x,y) \) 为两条垂直切线的交点。

\textbf{核心公式：}
\[
\boxed{x^{2}+y^{2}=a^{2}+b^{2}}
\]
该圆称为椭圆 \( \dfrac{x^{2}}{a^{2}}+\dfrac{y^{2}}{b^{2}}=1 \) 的蒙日圆。

\textbf{使用场景：} 已知两条切线垂直求交点轨迹；由交点在蒙日圆上推得对应两切线垂直；处理圆锥曲线切线垂直相关的定点、定值及最值问题。

\textbf{简单例子：} 椭圆 \( \dfrac{x^{2}}{9}+\dfrac{y^{2}}{4}=1 \) 中 \( a=3,\;b=2 \)，其蒙日圆为 \( x^{2}+y^{2}=13 \)。取点 \( P(3,2) \) 满足 \( 3^{2}+2^{2}=13 \)，由 \( P \) 向椭圆可引两条切线，且这两条切线互相垂直；反之，若椭圆外一点所引两切线垂直，则该点必在蒙日圆上。